\section{Introducción y descripción del problema}
\label{sec:introduction}

El dengue es una enfermedad viral transmitida por el mosquito \emph{Aedes aegypti}, que se reproduce en agua
acumulada. En el Perú es un problema de salud pública creciente: el sistema de vigilancia registra más de un
millón de casos entre 2000 y 2025, y las epidemias de 2023 y 2024 fueron las más grandes del periodo.

Este proyecto usa datos abiertos oficiales del Estado peruano para responder seis preguntas sobre el dengue:
cómo evolucionó, dónde y en quiénes tiene mayor incidencia, y cómo se relaciona con la cobertura de
establecimientos de salud, las emergencias hídricas y la lluvia. Todo el procesamiento se hace con las
Structured APIs de Apache Spark, la DataFrame API y Spark SQL \cite{zaharia2016,armbrust2015}.

Además de los resultados, el proyecto responde la pregunta central del curso: \emph{¿cómo transforma Spark
las operaciones de análisis en una ejecución distribuida?} Para ello se siguen consultas reales del análisis
desde el logical plan hasta los stages y tasks que se ejecutan, y se mide el shuffle que generan.
